% !TEX program = xelatex
% Lernplatt - by Can Siebert
% Kurs 22 (Bo 143) - Tag 12 - Lehrskript
% Revenue Streams, Multi-Channel & Plattform-Skalierung
%
% Self-contained xelatex source. Kein biblatex, keine externen Bilder.

\documentclass[11pt,a4paper,oneside]{article}

% --- Encoding & Sprache --------------------------------------------------
\usepackage{polyglossia}
\setdefaultlanguage{german}

% --- Geometrie -----------------------------------------------------------
\usepackage[a4paper,top=2cm,bottom=2cm,outer=2.5cm,inner=2.5cm,bindingoffset=1.5cm]{geometry}

% --- Fonts ---------------------------------------------------------------
\usepackage{fontspec}
\defaultfontfeatures{Ligatures=TeX}

\IfFontExistsTF{Charter}{%
  \setmainfont{Charter}[Numbers=OldStyle]%
}{%
  \IfFontExistsTF{Noto Serif}{%
    \setmainfont{Noto Serif}%
  }{%
    \setmainfont{Liberation Serif}%
  }%
}

\IfFontExistsTF{Source Sans 3}{%
  \setsansfont{Source Sans 3}%
}{%
  \IfFontExistsTF{Source Sans Pro}{%
    \setsansfont{Source Sans Pro}%
  }{%
    \setsansfont{Liberation Sans}%
  }%
}

\newcommand{\caveatfontfallback}{\itshape}
\IfFontExistsTF{Caveat}{%
  \newfontfamily\caveatfont{Caveat}[Scale=1.15]%
}{%
  \newcommand\caveatfont{\caveatfontfallback}%
}

\setmonofont{Liberation Mono}

\usepackage{microtype}
\linespread{1.15}

\usepackage{xcolor}
\definecolor{lpInk}{RGB}{20,28,38}
\definecolor{lpAccent}{RGB}{22,90,140}
\definecolor{lpAccentLight}{RGB}{210,228,240}
\definecolor{lpAmber}{RGB}{222,138,30}
\definecolor{lpAmberLight}{RGB}{255,243,217}
\definecolor{lpAmberBorder}{RGB}{196,118,18}
\definecolor{lpGray}{RGB}{96,104,114}
\definecolor{lpGrayLight}{RGB}{238,240,243}

\usepackage[most]{tcolorbox}
\tcbuselibrary{breakable,skins}

\newtcolorbox{eselsbox}[1][Eselsbrücke]{%
  enhanced, breakable,
  colback=lpAmberLight,
  colframe=lpAmberBorder,
  boxrule=0.6pt,
  arc=2pt,
  borderline={0.4pt}{0pt}{lpAmberBorder, dashed},
  left=10pt,right=10pt,top=8pt,bottom=8pt,
  fonttitle=\sffamily\bfseries\color{lpAmberBorder},
  coltitle=lpAmberBorder,
  title={#1},
  attach boxed title to top left={xshift=8pt,yshift=-8pt},
  boxed title style={size=small, colback=lpAmberLight, colframe=lpAmberBorder, boxrule=0.4pt, arc=1pt},
  top=14pt
}

\newtcolorbox{merkbox}[1][Merksatz]{%
  enhanced, breakable,
  colback=lpAccentLight,
  colframe=lpAccent,
  boxrule=0.6pt,
  arc=2pt,
  left=10pt,right=10pt,top=8pt,bottom=8pt,
  fonttitle=\sffamily\bfseries\color{white},
  coltitle=white,
  title={#1},
  attach boxed title to top left={xshift=8pt,yshift=-8pt},
  boxed title style={size=small, colback=lpAccent, colframe=lpAccent, boxrule=0.4pt, arc=1pt},
  top=14pt
}

\usepackage{enumitem}
\setlist{itemsep=2pt, topsep=4pt, parsep=0pt}
\setlist[itemize,1]{label={\textbullet},leftmargin=1.6em}
\setlist[itemize,2]{label={\textendash},leftmargin=1.4em}
\setlist[enumerate,1]{label=\arabic*., leftmargin=1.8em}

\usepackage{tabularx}
\usepackage{array}
\usepackage{booktabs}
\usepackage{longtable}

\usepackage{fancyhdr}
\usepackage{lastpage}
\pagestyle{fancy}
\fancyhf{}
\renewcommand{\headrulewidth}{0.3pt}
\renewcommand{\footrulewidth}{0pt}
\fancyhead[L]{\footnotesize\sffamily\color{lpGray}Lernplatt \textperiodcentered{} Kurs 22 \textperiodcentered{} Tag 12}
\fancyhead[R]{\footnotesize\sffamily\color{lpGray}Revenue Streams \& Skalierung}
\fancyfoot[L]{\footnotesize\sffamily\color{lpGray}\textcopyright{} Can Siebert}
\fancyfoot[R]{\footnotesize\sffamily\color{lpGray}\thepage{} / \pageref{LastPage}}

\fancypagestyle{plain}{%
  \fancyhf{}%
  \renewcommand{\headrulewidth}{0pt}%
  \fancyfoot[C]{\footnotesize\sffamily\color{lpGray}\thepage{} / \pageref{LastPage}}%
}

\usepackage[explicit]{titlesec}

\titleformat{\section}[block]
  {\sffamily\Large\bfseries\color{lpAccent}}
  {\thesection.}{0.6em}{#1}
  [{\vspace{2pt}\color{lpAccent}\titlerule[0.6pt]}]

\titleformat{\subsection}[block]
  {\sffamily\large\bfseries\color{lpInk}}
  {\thesubsection}{0.6em}{#1}

\titleformat{\subsubsection}[block]
  {\sffamily\normalsize\bfseries\color{lpInk}}
  {}{0pt}{#1}

\titlespacing*{\section}{0pt}{14pt}{8pt}
\titlespacing*{\subsection}{0pt}{10pt}{4pt}
\titlespacing*{\subsubsection}{0pt}{8pt}{2pt}

\usepackage[hidelinks,unicode]{hyperref}

\usepackage{tikz}
\usetikzlibrary{shapes.geometric,arrows.meta,positioning,calc,fit,backgrounds}

\renewcommand{\contentsname}{\sffamily\Large\bfseries\color{lpAccent}Inhalt}

\newcommand{\akzent}[1]{{\caveatfont\color{lpAmber}#1}}
\newcommand{\fachbegriff}[1]{\textbf{#1}}
\newcommand{\quelle}[1]{{\footnotesize\color{lpGray}(#1)}}

% =========================================================================
\begin{document}

% --- COVER ---------------------------------------------------------------
\begin{titlepage}
  \pagestyle{empty}
  \vspace*{1.5cm}
  \noindent{\sffamily\footnotesize\color{lpGray}LERNPLATT \textperiodcentered{} BY CAN SIEBERT}\\[2pt]
  \noindent{\color{lpAccent}\rule{\linewidth}{1.2pt}}\\[4pt]
  \noindent{\sffamily\footnotesize\color{lpGray}MODUL 6 \textperiodcentered{} TAG 12 \textperiodcentered{} KURS 22 (Bo 143) \textperiodcentered{} 15.\,Juni 2026}

  \vspace{3cm}
  {\sffamily\fontsize{30}{34}\selectfont\bfseries\color{lpInk}Revenue Streams,\\Multi-Channel\\\& Plattform-\\Skalierung\par}

  \vspace{0.7cm}
  {\sffamily\large\color{lpGray}Erlösquellen optimieren, Vertriebskanäle bündeln,\\Wachstum messbar und steuerbar gestalten.\par}

  \vspace{1.5cm}
  {\caveatfont\LARGE\color{lpAmber}Lehrskript -- Tag 12 von 16}\par

  \vfill

  \noindent\begin{minipage}{0.55\linewidth}
    {\sffamily\footnotesize\color{lpGray}Lernpfad:}\\
    {\sffamily\small\color{lpInk}Wissen \textrightarrow{} Anwendung \textrightarrow{} Management}\\[10pt]
    {\sffamily\footnotesize\color{lpGray}Bearbeitungsdauer:}\\
    {\sffamily\small\color{lpInk}1 Kurstag}
  \end{minipage}\hfill
  \begin{minipage}{0.4\linewidth}
    \raggedleft
    {\sffamily\footnotesize\color{lpGray}Dozent:}\\
    {\sffamily\small\color{lpInk}Can Siebert}\\[10pt]
    {\sffamily\footnotesize\color{lpGray}Format:}\\
    {\sffamily\small\color{lpInk}Theorie \textperiodcentered{} Fallstudie \textperiodcentered{} Coaching}
  \end{minipage}

  \vspace{0.5cm}
  \noindent{\color{lpAccent}\rule{\linewidth}{0.6pt}}
\end{titlepage}

% --- LERNZIELE -----------------------------------------------------------
\section*{Lernziele dieses Tages}
\addcontentsline{toc}{section}{Lernziele dieses Tages}
\thispagestyle{plain}

\noindent An Tag 11 haben Sie verstanden, wie eine Plattform überhaupt Erlöse generiert. Heute wird die Perspektive größer: Plattformgeschäfte verfügen typischerweise über mehrere Revenue Streams, mehrere Vertriebskanäle und mehrere Wachstumshebel \emph{gleichzeitig}. Die zentrale Frage lautet: Welche dieser Quellen werden \emph{skaliert}, welche \emph{stabilisiert}, welche bewusst \emph{begrenzt} -- und woran erkennt man im laufenden Betrieb, ob das Ganze funktioniert?

\subsection*{L1 -- Wissen (Reproduktion und Einordnung)}
\begin{itemize}
  \item Revenue Streams und Multi-Channel-Ansätze verstehen.
  \item Diversifikation von Erlösquellen einordnen.
  \item Performance-Analyse über KPIs, Conversion und Skalierungskennzahlen verstehen.
  \item Grundlagen von Wachstumsstrategien in Plattformmodellen kennen.
\end{itemize}

\subsection*{L2 -- Anwendung (Transfer auf konkrete Situationen)}
\begin{itemize}
  \item Bestehende Revenue Streams bewerten und Optimierungspotenziale erkennen.
  \item Multi-Channel-Ansätze für Plattformgeschäftsmodelle entwickeln.
  \item KPIs zur Steuerung von Umsatz, Conversion, Marge, Kundenwert und Wachstum anwenden.
  \item Conversion-Engpässe identifizieren und Wachstumsmaßnahmen ableiten.
\end{itemize}

\subsection*{L3 -- Management (Entscheidung und Verantwortung)}
\begin{itemize}
  \item Revenue-Architekturen strategisch gestalten und skalieren.
  \item Zielkonflikte zwischen Wachstum, Marge, Kanalabhängigkeit, Kundenwert und Komplexität bewerten.
  \item Entscheidungen unter Unsicherheit, Budgetrestriktionen und Zeitdruck treffen.
  \item Verantwortung für nachhaltige Skalierung und wirtschaftliche Steuerbarkeit übernehmen.
\end{itemize}

\begin{merkbox}[Roter Faden]
Mehrere Revenue Streams sind nicht automatisch ``mehr Geschäft''. Sie sind \emph{mehr Komplexität}. Erst eine geordnete \emph{Revenue-Architektur} -- mit klarer Rolle pro Quelle, eindeutigen KPIs und expliziter Skalierungs- bzw.\ Stabilisierungslogik -- verwandelt Diversifikation in Resilienz. Wer ohne Architektur diversifiziert, kauft sich Steuerungsverlust.
\end{merkbox}

\clearpage

% --- 1. EINSTIEG -------------------------------------------------------
\section{Einstieg: Vom Umsatz zur Architektur}

In vielen wachsenden Plattformen sehen die Berichte des Controllings ungefähr so aus: ``Umsatz pro Kanal X, pro Produktgruppe Y, pro Region Z''. Was selten klar wird, ist die \emph{Architektur} dahinter: Welche Quelle trägt die Marge, welche bringt das Volumen, welche die Differenzierung, welche das Risiko? Und welche zwei Quellen kannibalisieren sich vielleicht gerade unbemerkt? \quelle{Osterwalder \& Pigneur, 2010}

\subsection{Drei typische Wachstumspathologien}

\begin{enumerate}
  \item \fachbegriff{Ein-Bein-Geschäft.} 80\,\% des Umsatzes hängen an einer einzelnen Erlösquelle. Wachstum funktioniert, bis dieser eine Hebel kippt -- Algorithmus-Änderung, Provisionserhöhung, Wettbewerber-Einstieg.
  \item \fachbegriff{Diversifikations-Aktivismus.} Jeder neue Vorschlag wird gestartet -- ohne dass alte Quellen sauber gesteuert werden. Das Unternehmen wird hochaktiv, aber nicht profitabler.
  \item \fachbegriff{KPI-Inflation.} Es werden 80 Kennzahlen gemessen, aber niemand kann auf einer Folie zeigen, ob die Plattform gerade wirtschaftlich gewinnt oder verliert.
\end{enumerate}

\subsection{Was eine ``Revenue-Architektur'' ist}

Eine Revenue-Architektur ist mehr als eine Auflistung von Erlösquellen. Sie verbindet vier Ebenen:

\begin{itemize}
  \item \textbf{Quellen.} Welche Revenue Streams gibt es, in welcher Größe, mit welcher Marge?
  \item \textbf{Kanäle.} Über welche Vertriebskanäle werden sie erreicht? Eigene Plattform, Marktplatz, Partner, Direkt?
  \item \textbf{KPIs.} Woran wird Performance je Quelle / Kanal gemessen?
  \item \textbf{Steuerung.} Welche Quellen skalieren, welche stabilisieren, welche werden zurückgebaut?
\end{itemize}

\subsection{Senior-Level-Frage}

Die zentrale Senior-Level-Frage lautet nicht ``Wie können wir mehr Umsatz machen?'', sondern: \emph{Welche Revenue Streams sollen skaliert, welche stabilisiert und welche bewusst begrenzt werden?} -- und mit welcher Begründung. \quelle{Cusumano et al., 2019, Kap.~7}

\begin{eselsbox}[Eselsbrücke: \akzent{Skalieren -- Stabilisieren -- Stoppen}]
Drei S, die das Management einer reifen Plattform regelmäßig anwendet:\\[2pt]
\textbf{S}kalieren -- viel mehr investieren, schnell wachsen.\\
\textbf{S}tabilisieren -- halten und effizient betreiben.\\
\textbf{S}toppen / reduzieren -- bewusst beenden oder schrumpfen.\\[2pt]
\emph{Wer alle Quellen gleichzeitig ``ausbauen'' will, baut keine Architektur, sondern Aktivität.}
\end{eselsbox}

% --- 2. REVENUE STREAMS: PRIMÄR, SEKUNDÄR, STRATEGISCH ----------------
\section{Revenue Streams: primär, sekundär, strategisch}

Nicht alle Revenue Streams sind gleichberechtigt. Eine hilfreiche Drei-Teilung trennt:

\begin{enumerate}
  \item \fachbegriff{Primäre Revenue Streams.} Tragen einen wesentlichen Anteil des Umsatzes -- typischerweise 50--80\,\%. Sie definieren das Geschäftsmodell.
  \item \fachbegriff{Sekundäre Revenue Streams.} Ergänzen das Hauptgeschäft, glätten Schwankungen, erschließen Zusatzsegmente. Typischerweise 10--30\,\%.
  \item \fachbegriff{Strategische Revenue Streams.} Klein in Volumen, groß in Strategie -- sie bauen Optionen für die Zukunft (Datenservices, neue Marktsegmente, Pilotmodelle).
\end{enumerate}

\subsection{Typische Plattform-Revenue-Streams}

\begin{itemize}
  \item Transaktionsprovisionen -- meist primär.
  \item Subscription-Modelle (Anbieter / Endkunden) -- primär oder sekundär.
  \item Premium-Funktionen, Add-Ons -- sekundär.
  \item Werbung / Sponsored Listings -- sekundär (mit Vorsicht).
  \item Affiliate-Provisionen -- sekundär.
  \item Servicepakete (Onboarding, Beratung, Schulung) -- sekundär.
  \item Datenservices und Reporting-Produkte -- strategisch.
  \item Upgrades und Zusatzleistungen -- sekundär.
\end{itemize}

\subsection{Multi-Channel-Ansatz}

Ein \fachbegriff{Multi-Channel-Ansatz} bedeutet, dass dieselbe Plattform mehrere Vertriebskanäle nutzt, um ihre Revenue Streams zu erreichen: \quelle{Osterwalder \& Pigneur, 2010; Homburg, 2020}

\begin{itemize}
  \item Eigener Webshop / Plattform-Frontend.
  \item Marktplatz-Listing (extern).
  \item Partnervertrieb (Reseller, Distributoren).
  \item Direktvertrieb / Außendienst.
  \item Affiliate-Kanäle.
  \item Kundenportal (Bestandskunden, Wiederbestellung).
  \item E-Mail-Strecken (siehe Tag 10).
  \item Digitale Kampagnen (Suche, Social, Display).
\end{itemize}

\subsection{Diversifikation -- Vorteile und Risiken}

\begin{tabularx}{\linewidth}{@{}lXX@{}}
\toprule
& \textbf{Vorteile} & \textbf{Risiken} \\
\midrule
Mehrere Quellen & geringere Abhängigkeit, bessere Resilienz & höhere Komplexität \\
Mehrere Kanäle & breitere Reichweite, Kundenausschöpfung & Kannibalisierung, widersprüchliche Preise \\
Mehrere Segmente & größerer adressierbarer Markt & verstreute Ressourcen, unklare Positionierung \\
\bottomrule
\end{tabularx}

\begin{merkbox}[Diversifikations-Test]
Diversifikation lohnt sich, wenn \emph{zwei} Bedingungen erfüllt sind:
\begin{enumerate}
  \item Die neue Quelle hat eine \emph{nachweisbar andere} Risikologik als die bestehenden.
  \item Die Organisation hat \emph{eindeutige Verantwortliche} für die neue Quelle.
\end{enumerate}
Ist nur eines erfüllt, ist es keine Diversifikation, sondern Aktivismus.
\end{merkbox}

% --- 3. KANNIBALISIERUNG --------------------------------------------
\section{Kannibalisierung und Kanalkonflikte}

Mehrere Kanäle nebeneinander zu betreiben heißt fast immer: \emph{Kannibalisierung} ins Auge fassen. Ein Kunde, der bisher beim Außendienst kauft, kauft jetzt im Webshop -- formal ist es Umsatz, faktisch ist es ein Verschieben aus einem teureren in einen günstigeren Kanal.

\subsection{Vier typische Kanalkonflikte}

\begin{enumerate}
  \item \textbf{Preis.} Der Webshop ist günstiger als der Außendienst -- der Außendienst verliert Bestandskunden.
  \item \textbf{Sortiment.} Auf der externen Marktplatz-Plattform sind nicht alle Produkte verfügbar -- Frustration bei Wiederholungskäufern.
  \item \textbf{Information.} CRM des Außendienstes und Kundenportal sind nicht synchronisiert -- der Kunde wird im Außendienst nach Informationen gefragt, die er im Portal längst eingegeben hat.
  \item \textbf{Beziehung.} Plattformkunden ``gehören'' der Plattform -- der eigentliche Anbieter sieht den Kontakt nicht.
\end{enumerate}

\subsection{Kanalrollen klar definieren}

Eine wirksame Antwort auf Kanalkonflikte ist die \fachbegriff{Kanalrollen-Architektur}: \quelle{McAfee \& Brynjolfsson, 2017}

\begin{tabularx}{\linewidth}{@{}>{\bfseries}lX@{}}
\toprule
Kanal & Hauptfunktion \\
\midrule
Eigener Webshop & Bestellung von Bestandskunden, Premium-Beratung, Datenhoheit \\
Marktplatz / Plattform & Neukundengewinnung, Reichweite \\
Direktvertrieb / Außendienst & Großkunden, Beratungsprodukte, individuelle Angebote \\
Partner / Reseller & Nischen, Regionen, Branchenexpertise \\
Affiliate-Kanal & breite Lead-Generierung, Markenaufbau \\
Kundenportal & Wiederbestellung, Service, Self-Service \\
\bottomrule
\end{tabularx}

\medskip

\noindent\emph{Lesehinweis.} Jeder Kanal sollte eine eindeutige Hauptfunktion haben -- sonst wird er entweder selten genutzt oder konkurriert intern mit einem anderen.

% --- 4. PERFORMANCE-ANALYSE: KPIs --------------------------------------
\section{Performance-Analyse: KPIs entlang des Funnels}

Eine reife Plattform hat nicht ``viele KPIs'', sondern \emph{die richtigen} -- entlang eines klar definierten Funnels von Sichtbarkeit bis Umsatzwirkung.

\subsection{Der Plattform-Funnel}

\begin{enumerate}
  \item \fachbegriff{Sichtbarkeit.} Wie oft taucht ein Angebot in Suchanfragen / Empfehlungen auf?
  \item \fachbegriff{Klick.} Wie oft wird das Angebot tatsächlich angeklickt?
  \item \fachbegriff{Registrierung.} Wie viele Nutzer legen ein Konto an?
  \item \fachbegriff{Anfrage.} Wie viele stellen eine Anfrage oder buchen eine Demo?
  \item \fachbegriff{Kauf.} Wie viele schließen ab?
  \item \fachbegriff{Wiederkauf.} Wie viele bestellen erneut?
  \item \fachbegriff{Upgrade / Cross-Sell.} Wie viele wandern in ein höherwertiges Segment?
\end{enumerate}

\subsection{KPI-Cluster}

\begin{itemize}
  \item \textbf{Umsatz-Cluster.} Umsatz je Kanal, Deckungsbeitrag, Net Revenue Retention.
  \item \textbf{Conversion-Cluster.} Conversion Rate je Funnel-Stufe, Anteil Wiederkäufer, durchschnittlicher Warenkorb.
  \item \textbf{Kundenwert-Cluster.} CAC, CLV, Churn Rate, ARPU.
  \item \textbf{Wachstums-Cluster.} Kundenneuwachstum, Aktivierungsrate, Skalierungskennzahlen.
  \item \textbf{Partner-Cluster.} Partnerumsatz, Partner-CAC, Partner-Konzentration.
\end{itemize}

\subsection{Net Revenue Retention -- die unterschätzte Kennzahl}

Eine besonders wichtige Kennzahl in der Plattform- und Subscription-Welt ist die \fachbegriff{Net Revenue Retention (NRR)}: Wie viel Umsatz erzeugt eine Kunden-Kohorte heute, verglichen mit dem, was sie vor 12 Monaten erzeugt hat -- bei \emph{gleichem Kundenstamm}, also abzüglich Churn, aber inklusive Upsell. \quelle{Kaplan \& Norton, 1996; Cusumano et al., 2019}

\begin{itemize}
  \item NRR $<$ 100\,\% -- die Plattform verliert Umsatz pro Kohorte (Churn $>$ Upsell).
  \item NRR $\approx$ 100\,\% -- Wachstum kommt ausschließlich aus Neukunden.
  \item NRR $>$ 100\,\% -- die Plattform wächst \emph{innerhalb} der Bestandsbasis. Sehr starkes Signal.
\end{itemize}

\begin{merkbox}[Die Pyramide -- KPIs mit Schichten]
Operative Teams brauchen Aktivitäts-KPIs (Klicks, Anfragen). Mittleres Management braucht Conversion-KPIs. Geschäftsführung braucht Wirtschafts-KPIs (Deckungsbeitrag, NRR, CLV/CAC). Wenn alle drei Ebenen \emph{dieselben} Kennzahlen haben, fehlt die Hierarchie. Wenn keine Ebene die andere verstehen kann, fehlt die Brücke.
\end{merkbox}

% --- 5. CONVERSION-OPTIMIERUNG ---------------------------------------
\section{Conversion-Optimierung: weniger Traffic, mehr Wirkung}

Die häufigste Wachstumsempfehlung ``mehr Traffic'' ist meist die teuerste. \fachbegriff{Conversion-Optimierung} entlang des bestehenden Funnels ist in der Regel deutlich wirtschaftlicher. \quelle{Ellis \& Brown, 2017}

\subsection{Wo Conversion-Engpässe sitzen}

In der Praxis sitzen die größten Engpässe typischerweise an einer von vier Stellen:

\begin{enumerate}
  \item \textbf{Sichtbarkeit $\rightarrow$ Klick.} Die Plattform sieht das Angebot zwar, aber niemand klickt. Ursache: schwache Titel, schwache Vorschaubilder, hohe Konkurrenz im Ranking.
  \item \textbf{Klick $\rightarrow$ Anfrage.} Der Nutzer landet auf der Seite -- und verlässt sie wieder. Ursache: schlechte Produktseite, fehlende Vertrauenselemente, unklarer Preis.
  \item \textbf{Anfrage $\rightarrow$ Kauf.} Der Nutzer fragt an, kauft aber nicht. Ursache: langsame Reaktion, unklares Angebot, fehlende Beratung.
  \item \textbf{Erstkauf $\rightarrow$ Wiederkauf.} Der Kunde kommt nicht zurück. Ursache: schwache Onboarding-Strecke, mangelnde Bindung, kein Folgeangebot.
\end{enumerate}

\subsection{Vier Optimierungs-Werkzeuge}

\begin{itemize}
  \item \textbf{A/B-Tests.} Variante A vs.\ B in laufendem Traffic -- harte Datengrundlage.
  \item \textbf{Usability-Reviews.} Kleine Stichprobe echter Nutzer beobachtet beim Durchgang.
  \item \textbf{Funnel-Analyse.} Wo bricht der Trichter konkret ein?
  \item \textbf{Heatmaps / Session Recordings.} Was klicken Nutzer wirklich an?
\end{itemize}

\subsection{Ressourcenpriorisierung}

Conversion-Optimierung lohnt sich besonders dort, wo \emph{viel Traffic} auf \emph{wenig Abschluss} trifft -- das ist die Stelle mit der größten ökonomischen Hebelwirkung pro investiertem Euro. \quelle{McGrath, 2010}

\begin{eselsbox}[Eselsbrücke: \akzent{Hebel vor Volumen}]
Eine Conversion-Rate-Verbesserung von 1{,}5\,\% auf 2\,\% ist ökonomisch oft mehr wert als eine 50\,\%-ige Traffic-Erhöhung. \emph{Hebel vor Volumen.} Wer das versteht, verbrennt weniger Marketingbudget.
\end{eselsbox}

% --- 6. SKALIERUNGSSTRATEGIEN ----------------------------------------
\section{Skalierungsstrategien}

Skalierung ist die nächste Stufe nach Wachstum: planbare, wiederholbare Vervielfachung -- nicht ``mehr Aufwand für mehr Ergebnis'', sondern \emph{disproportionaler} Ergebniszuwachs.

\begin{figure}[h]
\centering
\begin{tikzpicture}[font=\sffamily\footnotesize,
  axis/.style={->, >=Stealth, draw=lpGray, line width=0.5pt},
  curveW/.style={draw=lpAccent, line width=1pt},
  curveS/.style={draw=lpAmberBorder, line width=1pt},
  lab/.style={font=\sffamily\scriptsize, text=lpInk}]
  \draw[axis] (0,0) -- (7.4,0) node[right, lab] {Aufwand};
  \draw[axis] (0,0) -- (0,4) node[above, lab] {Ergebnis};
  \draw[curveW] (0,0) -- (7,3.2);
  \node[text=lpAccent, font=\sffamily\small\bfseries] at (6.0,2.3) {Wachstum};
  \node[text=lpAccent, font=\sffamily\scriptsize, align=center] at (6.0,1.7) {linear:\\Aufwand $\to$ Ergebnis};
  \draw[curveS] plot[domain=0:7, samples=40, smooth] (\x, {0.18*\x*\x});
  \node[text=lpAmberBorder, font=\sffamily\small\bfseries] at (5.0,3.6) {Skalierung};
  \node[text=lpAmberBorder, font=\sffamily\scriptsize, align=center] at (5.0,3.05) {disproportional:\\Hebel statt Aufwand};
  \draw[draw=lpGray, dashed, line width=0.3pt] (3,0) -- (3,3.2);
  \node[lab, below, text=lpGray] at (3,0) {Stabilisierungs-Schwelle};
\end{tikzpicture}
\caption{\sffamily Wachstum ist linear -- mehr Aufwand bringt mehr Ergebnis. Skalierung erzeugt einen disproportionalen Effekt; voraussetzungsreich, aber stark hebelnd.}
\label{fig:scale-22t12}
\end{figure}

\subsection{Sechs Skalierungsachsen}

\begin{enumerate}
  \item \fachbegriff{Kanal-Skalierung.} Mehr Volumen über bestehende Kanäle (Plattform, Webshop, Partner).
  \item \fachbegriff{Kundensegment-Skalierung.} Erschließung neuer Zielgruppen (z.\,B.\ vom Mittelstand zum Konzern, oder umgekehrt).
  \item \fachbegriff{Geografische Skalierung.} Expansion in neue Märkte / Länder / Regionen.
  \item \fachbegriff{Produkt- und Serviceausbau.} Neue Funktionen, Module, Service-Linien.
  \item \fachbegriff{Partnernetzwerke.} Mehr Partner, neue Partnertypen, exklusive Allianzen.
  \item \fachbegriff{Automatisierung und Daten.} Datengetriebene Optimierung, KI-Unterstützung, automatisierte Prozesse.
\end{enumerate}

\subsection{Eine Achse zur Zeit}

Eine wichtige Senior-Level-Erkenntnis: \emph{Skalierung gleichzeitig auf mehreren Achsen ist riskant.} Wer parallel in zwei neue Länder \emph{und} ein neues Kundensegment \emph{und} eine neue Produktlinie expandiert, baut Komplexität exponentiell auf. Die meisten erfolgreichen Skalierungen bearbeiten eine Achse intensiv, stabilisieren, und gehen dann zur nächsten. \quelle{Cusumano et al., 2019, Kap.~7--8}

\subsection{Wachstumstypen unterscheiden}

\begin{itemize}
  \item \textbf{Organisches Wachstum} -- aus eigenen Aktivitäten, eigenem Cashflow.
  \item \textbf{Bezahltes Wachstum} -- über Werbung, Marketing, externe Lead-Generierung.
  \item \textbf{Virales Wachstum} -- Nutzer ziehen weitere Nutzer nach (Empfehlung, Netzwerkeffekt).
  \item \textbf{Partner-Wachstum} -- über Partnernetzwerke.
  \item \textbf{Akquisitorisches Wachstum} -- M\&A, Übernahmen.
\end{itemize}

% --- 7. WACHSTUMS-ENTSCHEIDUNGEN ------------------------------------
\section{Wachstums- und Reduktions-Entscheidungen}

Senior-Level-Steuerung erkennt nicht nur, was \emph{wachsen} soll, sondern auch, was \emph{schrumpfen} oder \emph{enden} soll. Diese Entscheidung wird in der Praxis vermieden -- mit teuren Folgen.

\subsection{Vier Entscheidungstypen pro Revenue Stream}

\begin{enumerate}
  \item \textbf{Skalieren.} Wesentliche Investition, eindeutiger Wachstumsplan, klare KPIs.
  \item \textbf{Testen.} Pilotbudget, definierte Testdauer, klare Ja/Nein-Kriterien am Ende.
  \item \textbf{Stabilisieren.} Halten, effizient betreiben, nicht weiter ausbauen.
  \item \textbf{Reduzieren / Beenden.} Bewusst zurückbauen oder ganz schließen, Ressourcen umlenken.
\end{enumerate}

\subsection{Wann eine Quelle bewusst beendet wird}

\begin{itemize}
  \item Wenn der Deckungsbeitrag negativ ist und keine Aussicht auf Verbesserung besteht.
  \item Wenn die Quelle organisatorische Ressourcen bindet, die wertvoller eingesetzt werden könnten.
  \item Wenn die Quelle Reputations- oder Compliance-Risiken erzeugt, die nicht zur Strategie passen.
  \item Wenn die Quelle ein anderes (wichtigeres) Geschäft kannibalisiert.
\end{itemize}

\subsection{Beispiel: Portfolio-Entscheidung}

\begin{tabularx}{\linewidth}{@{}>{\bfseries}l X l l@{}}
\toprule
Revenue Stream & Status / Argument & Entscheidung & Budget \\
\midrule
Transaktionsprovision & Hauptträger, stabil & Stabilisieren & Halten \\
Subscription Enterprise & Wächst stark, CLV $\gg$ CAC & Skalieren & Stark erhöhen \\
Premium-Profil & Schwache Adoption, unklar & Testen 6\,M & Pilotbudget \\
Werbeplätze & Niedrige Marge, Reibung & Beenden & 0 \\
Datenservices & Strategisch interessant & Testen 12\,M & Klein \\
\bottomrule
\end{tabularx}

% --- 7b. ORGANISATIONSREIFE -----------------------------------
\section{Organisationsreife und Verantwortlichkeiten}

Wachstum bricht in Plattformen meist nicht an der Strategie, sondern an der Organisation. Eine schnell wachsende Revenue-Architektur braucht eine mitwachsende Organisationsreife.

\subsection{Vier kritische Verantwortungsfelder}

\begin{enumerate}
  \item \fachbegriff{Revenue Ownership.} Pro Stream eine klare Person mit Umsatz- und Margenverantwortung.
  \item \fachbegriff{Channel Ownership.} Pro Kanal eine klare Person, die die Kanalrolle steuert.
  \item \fachbegriff{Data Ownership.} Pro Datenfeld eine Person, die für Qualität, Aktualität und Definition steht.
  \item \fachbegriff{Customer Ownership.} Pro Segment ein klar verantwortlicher Account.
\end{enumerate}

\subsection{Typische Reibungspunkte}

\begin{itemize}
  \item Marketing und Plattformmanagement streiten über Lead-Qualität.
  \item Vertrieb und Account Management streiten über Bestandskunden.
  \item Partner Management und Direktvertrieb streiten über Provisionen.
  \item Controlling sieht andere Zahlen als die operativen Bereiche.
\end{itemize}

\subsection{Governance-Strukturen}

Eine reife Revenue-Governance hat vier Bausteine:

\begin{itemize}
  \item \textbf{Revenue Council.} Monatliches Gremium, das Revenue Streams, Kanäle und KPIs gemeinsam reviewt.
  \item \textbf{Channel Conflict Resolution.} Klares Verfahren bei Kanalkonflikten (Preis, Sortiment, Bestandskunden).
  \item \textbf{KPI-Definitionen.} Verbindliche Definitionen pro Kennzahl -- ``unsere Conversion Rate'' bedeutet überall dasselbe.
  \item \textbf{Investitions-Logik.} Klare Regeln, wer wann Budget für welche Quelle bekommt.
\end{itemize}

\begin{merkbox}[Organisations-Faustregel]
Wenn die Vertriebsorganisation länger braucht, um eine Quartalsauswertung zu konsolidieren, als das Quartal lang ist, fehlt nicht das Reporting -- es fehlt die Organisationsreife. Die nächste Skalierungsstufe ist dann \emph{nicht} ein weiterer Revenue Stream, sondern eine bessere Governance.
\end{merkbox}

% --- 8. MANAGEMENT-PERSPEKTIVE --------------------------------------
\section{Management-Perspektive: Zielkonflikte}

Wie schon an anderen Tagen gilt: die wichtigsten Entscheidungen sind nicht zu \emph{lösen}, sondern bewusst zu \emph{entscheiden}.

\subsection{Fünf zentrale Zielkonflikte}

\begin{enumerate}
  \item \textbf{Wachstum vs.\ Marge.} Schnelles Wachstum kostet Marge -- durch Werbung, Rabatte, neue Hilfsstrukturen.
  \item \textbf{Kanalabhängigkeit vs.\ Reichweite.} Ein dominanter Kanal bringt Reichweite und ein systemisches Risiko.
  \item \textbf{Kundenwert vs.\ Komplexität.} Mehrere Tiers, mehrere Kanäle und mehrere Subscription-Modelle erhöhen den durchschnittlichen Kundenwert -- und die operative Komplexität.
  \item \textbf{Skalierung vs.\ Steuerbarkeit.} Je schneller skaliert wird, desto schwerer wird Steuerung. Datenqualität, Reporting, Prozessdisziplin müssen mitwachsen.
  \item \textbf{Wachstum vs.\ Resilienz.} Eine Plattform, die nur wächst, ohne Diversifikation, ist verletzlich gegen den einen Schock, der trifft.
\end{enumerate}

\subsection{Senior-Level-Entscheidungslogik}

\begin{itemize}
  \item \textbf{Trenne Mengen-KPIs von Wert-KPIs.} ``Mehr Anfragen'' ist nicht ``mehr Geschäft''.
  \item \textbf{Erkenne Kannibalisierung früh.} Ein Kanal-Erfolg auf Kosten eines anderen ist kein Wachstum.
  \item \textbf{Erlaube Beendigungen.} ``Mehr Quellen'' ist nicht automatisch besser. Manche Quellen sollte man schließen.
  \item \textbf{Investiere in NRR.} Wachstum aus dem Bestand ist nachhaltiger als Wachstum aus Neukunden.
\end{itemize}

\subsection{Die Komplexitätsgrenze}

Es gibt eine empirische Grenze, ab der zusätzliche Diversifikation mehr Komplexitätskosten als Resilienzgewinn erzeugt. Diese Grenze ist organisations-spezifisch -- sie hängt von Reife des Controllings, Datenqualität, Führungstiefe und Prozessen ab. Die Aufgabe des Senior-Level-Managements ist, diese Grenze zu kennen und bei Wachstumsplänen \emph{vor} Überschreitung zu adressieren -- nicht erst, wenn das Reporting zerfällt. \quelle{McGrath, 2010; Cusumano et al., 2019}

\begin{eselsbox}[Eselsbrücke: \akzent{W-M-K-S-R}]
Die fünf Zielkonflikte als Reihenfolge:\\[2pt]
\textbf{W}achstum vs.\ Marge -- \textbf{M}aturity (Kanalabhängigkeit) -- \textbf{K}omplexität -- \textbf{S}teuerbarkeit -- \textbf{R}esilienz.\\[2pt]
\emph{Mehr Quellen sind nur dann ein Gewinn, wenn die fünf W-M-K-S-R-Punkte abgewogen sind.}
\end{eselsbox}

\begin{merkbox}[Schluss-Merksatz für Tag 12]
Eine reife Revenue-Architektur ist nicht \emph{die Summe} aller möglichen Erlösquellen. Sie ist eine \emph{Auswahl} -- bewusst getroffen, mit klarer Rolle pro Quelle, eindeutigen KPIs und expliziter Skalierungs-, Stabilisierungs- oder Beendigungsentscheidung. Das ist der Unterschied zwischen ``viel Umsatz'' und ``steuerbarem Geschäft''.
\end{merkbox}

% --- 9. LITERATUR -----------------------------------------------------
\clearpage
\section{Literatur}

Im Skript verwendete und für die Vertiefung empfohlene Quellen. Zitierstil: APA-7.

\begin{description}[style=nextline,leftmargin=18pt,labelindent=0pt,itemsep=4pt]
  \item[Cusumano, M.\,A., Gawer, A., \& Yoffie, D.\,B. (2019).] \emph{The business of platforms: Strategy in the age of digital competition, innovation, and power}. HarperBusiness.
  \item[Ellis, S., \& Brown, M. (2017).] \emph{Hacking growth: How today's fastest-growing companies drive breakout success}. Crown Business.
  \item[Homburg, C. (2020).] \emph{Marketingmanagement: Strategie -- Instrumente -- Umsetzung -- Unternehmensführung} (7.\ Aufl.). Springer Gabler.
  \item[Iyer, B., \& Singh, H. (2018).] Designing platforms for value capture. \emph{MIT Sloan Management Review, 60}(1), 1--10.
  \item[Kaplan, R.\,S., \& Norton, D.\,P. (1996).] \emph{The balanced scorecard: Translating strategy into action}. Harvard Business School Press.
  \item[McAfee, A., \& Brynjolfsson, E. (2017).] \emph{Machine, platform, crowd: Harnessing our digital future}. W.\,W.\ Norton.
  \item[McGrath, R.\,G. (2010).] Business models: A discovery driven approach. \emph{Long Range Planning, 43}(2--3), 247--261. \href{https://doi.org/10.1016/j.lrp.2009.07.005}{https://doi.org/10.1016/j.lrp.2009.07.005}
  \item[Osterwalder, A., \& Pigneur, Y. (2010).] \emph{Business model generation: A handbook for visionaries, game changers, and challengers}. Wiley.
\end{description}

% --- 10. GLOSSAR ------------------------------------------------------
\clearpage
\section{Glossar}

Die wichtigsten Begriffe des Tages -- kurz definiert, in alphabetischer Reihenfolge.

\begin{description}[style=nextline,leftmargin=0pt,labelindent=0pt,itemsep=4pt]
  \item[A/B-Test] Vergleich zweier Varianten im laufenden Traffic zur datengetriebenen Optimierung.
  \item[ARPU] Average Revenue per User -- durchschnittlicher Umsatz pro Konto.
  \item[CAC] Customer Acquisition Cost -- Kosten zur Gewinnung eines zahlenden Kunden.
  \item[CLV] Customer Lifetime Value -- Wertsumme eines Kunden über die gesamte Beziehung.
  \item[Conversion Rate] Anteil der Besucher / Nutzer, die in die nächste Funnel-Stufe übergehen.
  \item[Churn Rate] Anteil der Kunden, die in einem Zeitraum kündigen.
  \item[Direktvertrieb] Vertrieb über eigenen Außendienst, ohne zwischengeschaltete Partner.
  \item[Funnel] Konversionspfad von Sichtbarkeit über Klick und Anfrage bis zum Kauf.
  \item[Heatmap] Visualisierung von Klick- und Aufmerksamkeitsverteilung auf einer Seite.
  \item[Kannibalisierung] Wachstum eines Kanals auf Kosten eines anderen.
  \item[Kanalrollen-Architektur] Klare Definition der Hauptfunktion pro Vertriebskanal.
  \item[Multi-Channel-Ansatz] Nutzung mehrerer paralleler Vertriebskanäle.
  \item[Net Revenue Retention (NRR)] Umsatzentwicklung einer Kunden-Kohorte (Bestand) inkl.\ Upsell, abzüglich Churn.
  \item[Partner-CAC] Akquisekosten pro Kunde, der über einen Partner gewonnen wurde.
  \item[Performance-Analyse] Strukturierte Bewertung von Aktivitäts-, Conversion- und Wirtschafts-KPIs.
  \item[Primäre Revenue Streams] Erlösquellen mit dominantem Umsatzanteil; tragen das Geschäftsmodell.
  \item[Reseller] Wiederverkäufer, der eigenständig an Endkunden verkauft.
  \item[Revenue-Architektur] Komposition mehrerer Erlösquellen mit klaren Rollen, KPIs und Steuerungsentscheidungen.
  \item[Revenue Stream] Einzelne Erlösquelle eines Geschäftsmodells.
  \item[Sekundäre Revenue Streams] Erlösquellen mit ergänzendem Charakter (10--30\,\% Anteil).
  \item[Skalierung] Disproportionale Vervielfachung von Output relativ zu Input.
  \item[Strategische Revenue Streams] Klein im Volumen, groß in der Optionalität für die Zukunft.
  \item[Upselling] Wechsel eines Bestandskunden in ein höheres Tier oder zu höherwertigem Angebot.
  \item[Virales Wachstum] Wachstum durch Empfehlung, bei dem bestehende Nutzer neue Nutzer mitbringen.
  \item[Wachstumstypen] Organisches, bezahltes, virales, partnerbasiertes, akquisitorisches Wachstum.
  \item[Zielkonflikte (Revenue / Skalierung)] Wachstum vs.\ Marge, Kanalabhängigkeit, Komplexität, Steuerbarkeit, Resilienz.
\end{description}

\vspace{1cm}
\noindent{\color{lpAccent}\rule{\linewidth}{0.6pt}}\\[4pt]
{\footnotesize\sffamily\color{lpGray}Lernplatt \textperiodcentered{} by Can Siebert \textperiodcentered{} Kurs 22 (Bo 143) \textperiodcentered{} Tag 12 \textperiodcentered{} \textcopyright{} 2026}

\end{document}
